\label{sec:introduction}

The geography of American voters is highly sorted. Urban cores are heavily Democratic,
exurban and rural areas are heavily Republican, and a broad suburban ring---expanding
and shifting through the 2010s and 2020s---is genuinely competitive~\citep{rodden2019why}.
This geographic sorting creates an inevitable tension in redistricting: any map that draws
compact, geographically coherent districts will reproduce some of this sorting at the
district level, producing safe urban Democratic seats and safe rural Republican seats.
The question is whether enacted maps produce \textit{more} safe seats than geography
requires---and if so, by how much.

This paper addresses that question by comparing safe-seat counts under \textsc{bisect}
algorithmic redistricting to enacted 2022 congressional maps in three states. The
\textsc{bisect} algorithm~\citep{bisect2024} draws districts by recursively bisecting
a census tract graph to optimise population balance, without reference to precinct-level
partisan data. The maps it produces represent a geographic lower bound on safe-seat
creation: they contain the safe seats that geographic partisan sorting makes unavoidable,
without the additional safe seats that enacted mapmakers create through deliberate packing
and cracking of partisan voter concentrations.

\subsection{Why Safe Seats Matter}

Safe seats have consequences for democratic representation that extend beyond the
incumbent's electoral security. \citet{fiorina1980incumbency} documented that safe-seat
incumbents are less responsive to constituent opinion than competitive-district members;
\citet{stonecash2003parties} showed that safe-seat representatives vote more ideologically
extreme than their districts' median voter. The entrenchment of safe seats across
redistricting cycles produces a Congress that is systematically less responsive to
statewide partisan shifts than the underlying vote distribution would imply.

From a legal perspective, states with ``free elections'' or ``fair elections'' clauses
in their constitutions have increasingly scrutinised safe-seat creation as a mechanism
of incumbent entrenchment. In \textit{League of Women Voters v.\ Commonwealth}\footnote{
\textit{League of Women Voters v.\ Commonwealth}, 178 A.3d 737 (Pa.\ 2018).} and
\textit{Harper v.\ Hall}\footnote{\textit{Harper v.\ Hall}, 868 S.E.2d 499 (N.C.\ 2022).}
the relevant courts held that redistricting plans that systematically reduce electoral
competition violate state constitutional free-elections guarantees. The \textsc{bisect}
comparison provides the critical counterfactual: how many competitive districts could
a compact, population-balanced map produce, and how does that compare to what the enacted
map produced?

\subsection{Paper Organisation}

Section~\ref{sec:background} defines safe seats and describes the measurement approach.
Section~\ref{sec:methodology} explains how \textsc{bisect} maps interact with the
geographic partisan sorting that produces safe seats. Section~\ref{sec:results}
presents the empirical results for NC, WI, and TX. Section~\ref{sec:discussion}
provides the party-specific decomposition and identifies which incumbents benefit from
enacted safe-seat creation. Section~\ref{sec:conclusion} concludes.
